% Part C -- Knowledge quiz, FORM A (20 items).
%
% Every item here has a counterpart in part_c_form_b.tex occupying the same
% blueprint cell (same topic, Bloom level, target difficulty and targeted
% misconception) with a different scenario and no shared stem wording.
% See docs/blueprint.md for the full mapping.
%
% Key positions are balanced: 5 items each with the key at (a), (b), (c), (d).
%
% The quiz was cut from 28 to 20 items on 2026-08-05. The eight retired A/B
% pairs are archived verbatim in common/deprecated_items.tex, which also records
% the old-to-new numbering.

\partheading{C}{Knowledge questions}

\instruction{20 questions, each with exactly one best answer. Tick one box per question. If you do not
know, please tick ``I do not know'' rather than guessing --- an honest blank is more useful to us than a
lucky guess, and none of this affects your grade. Plan for about 20 minutes.}

\begin{questions}

% ===== The model and its context =====================================

% --- Item 1: T1 / remember / easy -------------------------------------
\qitem The LLM in an agentic system is usually described as \emph{stateless across calls}. What does
that mean?
\begin{choices}
  \choice Its weights are frozen, so it cannot be adapted to new data.
  \CorrectChoice It retains nothing from one call to the next, so everything it needs must be supplied in the current request.
  \choice The provider stores the conversation and reloads it automatically for each call.
  \choice It returns the same output for the same input, because it holds no random state.
\end{choices}
\dontknow
\itemmeta{T1 --- LLM role \& statelessness}{remember}{easy}{A01}
\rationale{Statelessness is about the call boundary, not about training. (a) confuses statelessness with
frozen weights; (c) is the opposite of statelessness; (d) confuses it with determinism.}

% --- Item 2: T2 / understand / medium ---------------------------------
\qitem Why does a harness try to keep the context small even in runs that are nowhere near filling the
context window?
\begin{choices}
  \CorrectChoice Because model performance degrades as the window fills, so the goal is the smallest set of high-signal tokens that still produces the desired behaviour.
  \choice Because tokens that are not needed in the current step are discarded by the model anyway.
  \choice Because the context window is shared between all agents using the same endpoint.
  \choice Because a shorter context guarantees that the loop terminates in fewer iterations.
\end{choices}
\dontknow
\itemmeta{T2 --- Context engineering}{understand}{medium}{A02}
\rationale{This is context rot: the constraint is degradation, not the hard limit. Distractors treat the
window as a hard limit only, as shared state, or as a loop bound.}

% --- Item 3: T2 / apply-analyse / hard --------------------------------
\qitem A run has grown so long that the context no longer fits. The team wants to keep the agent
coherent while spending fewer tokens. Which pair of measures matches this goal?
\begin{choices}
  \choice Raise the temperature so the model relies less on the context, and shorten the system prompt.
  \choice Split the context evenly across several LLM calls and merge the answers afterwards.
  \CorrectChoice Summarise the older turns, and move state that is not needed right now to external storage that can be read back on demand.
  \choice Drop the tool schemas from the request, since the model has already used those tools once.
\end{choices}
\dontknow
\itemmeta{T2 --- Context engineering}{apply-analyse}{hard}{A03}
\rationale{Compaction plus offloading. (d) is attractive but breaks tool calling, because schemas are not
remembered across calls --- it also tests T1.}

% ===== The loop, tools, MCP and skills ===============================

% --- Item 4: T3 / understand / medium ---------------------------------
\qitem Why does the loop need guard conditions such as an iteration cap or a token budget?
\begin{choices}
  \CorrectChoice Because the LLM may return neither a final answer nor a valid action, and without a bound the loop would never terminate.
  \choice Because the model provider charges a penalty for requests above a fixed rate.
  \choice Because every iteration adds a tool to the registry, which would eventually exhaust the schema limit.
  \choice Because the user must be given a chance to intervene after a fixed number of steps.
\end{choices}
\dontknow
\itemmeta{T3 --- Loop and control flow}{understand}{medium}{A04}
\rationale{Guards guarantee termination. (d) confuses guard conditions with permission gates.}

% --- Item 5: T3 / apply-analyse / hard --------------------------------
\qitem A team extends their loop so that the agent first writes a plan for the task and then executes
it. What must the loop now do that a plain ReAct loop does not?
\begin{choices}
  \choice Call the LLM twice per iteration, once for reasoning and once for acting.
  \choice Persist the plan to long-term memory before any tool may be called.
  \choice Disable its iteration cap, because the plan already bounds the number of steps.
  \CorrectChoice Contrast the progress made so far against that initial plan on each pass through the loop.
\end{choices}
\dontknow
\itemmeta{T3 --- Loop and control flow}{apply-analyse}{hard}{A05}
\rationale{Up-front planning adds a comparison against the plan. (c) is the dangerous misconception:
a plan is not a termination guarantee.}

% --- Item 6: T4 / remember / easy -------------------------------------
\qitem Which of the following best describes a \emph{tool} in an agentic system?
\begin{choices}
  \choice A store of facts that survives from one agent run to the next.
  \CorrectChoice An interface with a defined schema through which the agent retrieves information or acts on an environment.
  \choice A prompt template that instructs the model how to solve a class of tasks.
  \choice The component that decides which action the agent should take next.
\end{choices}
\dontknow
\itemmeta{T4 --- Tools and tool design}{remember}{easy}{A06}
\rationale{The distractors are the neighbouring components: long-term memory (a), a skill (c), the LLM/loop (d).}

% --- Item 7: T4 / apply-analyse / hard --------------------------------
\qitem A \texttt{read\_file} tool returns an empty string both when the file is empty and when the path
does not exist. What is the most likely consequence for the agent?
\begin{choices}
  \choice It retries the call with the same arguments until the loop's iteration cap is reached.
  \choice The tool registry rejects the call, because an empty result violates the tool's schema.
  \CorrectChoice It cannot tell the two situations apart and reasons on from a wrong observation, because the failure was never communicated to it.
  \choice The loop enters its terminal state, because an empty observation carries no information.
\end{choices}
\dontknow
\itemmeta{T4 --- Tools and tool design}{apply-analyse}{hard}{A07}
\rationale{Tool failures must be communicated upstream. Silent, ambiguous success values are worse than
errors because the agent has no signal that anything went wrong.}

% --- Item 8: T5 / remember / easy -------------------------------------
\qitem What is the Model Context Protocol (MCP)?
\begin{choices}
  \choice A language model specialised in managing an agent's context window.
  \CorrectChoice A protocol through which an agent connects to servers that expose tools and other capabilities.
  \choice A file format for storing an agent's long-term memory between runs.
  \choice A specification of how much context a model must be able to process.
\end{choices}
\dontknow
\itemmeta{T5 --- MCP}{remember}{easy}{A08}
\rationale{(a) and (d) target the widespread ``MCP is a model / a context size standard'' confusion
invited by the name.}

% --- Item 9: T5 / apply-analyse / medium ------------------------------
\qitem A requirement says the tools of an MCP server must be made available to the model \emph{without
hard-coding individual tools}. What does the harness have to do?
\begin{choices}
  \choice Generate a Python wrapper function per tool before the application is built.
  \choice Keep a configuration file that lists every tool name the server is known to provide.
  \choice Ask the model which tools it thinks exist and validate its guesses against the server.
  \CorrectChoice Query the server for its tool list at runtime and pass the returned schemas on to the model.
\end{choices}
\dontknow
\itemmeta{T5 --- MCP}{apply-analyse}{medium}{A09}
\rationale{Runtime discovery. (a) and (b) are both hard-coding, just at different times.}

% --- Item 10: T6 / understand / medium --------------------------------
\qitem Why does a harness usually \emph{not} keep all available skills in the context at all times?
\begin{choices}
  \choice Because skills may contradict one another if more than one is loaded at a time.
  \choice Because the model can execute only one skill per run.
  \choice Because a skill has to be validated by a hook before it may enter the context.
  \CorrectChoice Because the context should carry only what the current task needs; a skill is disclosed when the task calls for it.
\end{choices}
\dontknow
\itemmeta{T6 --- Skills}{understand}{medium}{A10}
\rationale{Just-in-time disclosure follows from the context budget, not from conflicts between skills.}

% ===== Memory, hooks, permissions, sandboxing and sub-agents =========

% --- Item 11: T7 / remember / easy ------------------------------------
\qitem What characterises long-term memory, as opposed to context?
\begin{choices}
  \choice It is the part of the context window that is never summarised away.
  \choice It is the conversation history of the current run, kept in full.
  \CorrectChoice It lives outside the context window, persists across runs, and is consulted on demand.
  \choice It is the larger context window offered by models that support one.
\end{choices}
\dontknow
\itemmeta{T7 --- Long-term memory}{remember}{easy}{A11}
\rationale{(d) is the ``long-term memory is just a bigger window'' misconception.}

% --- Item 12: T8 / remember / easy ------------------------------------
\qitem What is a hook in an agent harness?
\begin{choices}
  \choice A tool the agent may call to inspect its own state.
  \CorrectChoice A custom function that runs at a defined point in the agent's lifecycle, such as before a tool call.
  \choice A connection to an external server from which tools are discovered.
  \choice A rule that decides whether the agent needs the user's approval.
\end{choices}
\dontknow
\itemmeta{T8 --- Hooks}{remember}{easy}{A12}
\rationale{Distractors are, in order, a tool, an MCP connection and a permission policy.}

% --- Item 13: T8 / apply-analyse / medium -----------------------------
\qitem A team wants to block any tool call whose arguments point outside the sandbox root, \emph{before}
it is executed. Which mechanism is designed for that?
\begin{choices}
  \CorrectChoice A pre-use hook that inspects the arguments and can reject the call.
  \choice A post-use hook that reports the violation once the tool has run.
  \choice A sub-agent that reviews the main agent's proposed actions.
  \choice A longer system prompt instructing the model never to leave the sandbox.
\end{choices}
\dontknow
\itemmeta{T8 --- Hooks}{apply-analyse}{medium}{A13}
\rationale{(d) is the ``prompt instead of enforce'' misconception; (b) has the right mechanism but the
wrong lifecycle point.}

% --- Item 14: T9 / understand / easy ----------------------------------
\qitem In a tiered permission model, which operation belongs in the tier that requires mandatory
confirmation by the user?
\begin{choices}
  \choice Listing the files in the agent's working directory.
  \choice Running the project's unit tests inside the container.
  \choice Reading a configuration file in the sandbox.
  \CorrectChoice Opening a pull request against a public repository.
\end{choices}
\dontknow
\itemmeta{T9 --- Permission management}{understand}{easy}{A14}
\rationale{The criterion is whether the consequences are irreversible or leave the sandbox, not how
complicated the operation is.}

% --- Item 15: T9 / understand / medium --------------------------------
\qitem A permission policy supports the outcomes \emph{allow}, \emph{deny} and \emph{require user
confirmation}. Why is the third outcome needed rather than deciding every case in advance?
\begin{choices}
  \choice Because the policy engine cannot evaluate its rules fast enough to keep the loop responsive.
  \choice Because the model cannot produce arguments that a static rule would be able to check.
  \CorrectChoice Because for some actions the consequences depend on the situation, so the judgement has to be handed to a human at the moment it arises.
  \choice Because denying an action would end the run, which the harness has to avoid.
\end{choices}
\dontknow
\itemmeta{T9 --- Permission management}{understand}{medium}{A15}
\rationale{Human-in-the-loop gates exist for context-dependent risk. (d) misstates what a denial does.}

% --- Item 16: T10 / understand / medium -------------------------------
\qitem Besides containment, what second benefit does a stable sandbox provide?
\begin{choices}
  \choice Lower latency, because the sandboxed environment sits closer to the model endpoint.
  \choice A guarantee that the agent cannot loop indefinitely.
  \choice Automatic approval of low-risk actions without asking the user.
  \CorrectChoice Reproducibility: the same command, seed and dependency set can be replayed under comparable conditions, so a failure reflects a real defect rather than environment drift.
\end{choices}
\dontknow
\itemmeta{T10 --- Sandboxing}{understand}{medium}{A16}
\rationale{Reproducibility is the second stated purpose of containerisation in this lab. (b) and (c)
attribute loop control and permission decisions to the sandbox.}

% --- Item 17: T11 / understand / medium -------------------------------
\qitem A main agent delegates a subtask to a specialised sub-agent. What does the main agent get back,
and why does that matter?
\begin{choices}
  \CorrectChoice A condensed result, so that the sub-agent's detailed intermediate work never enters the main agent's context.
  \choice The sub-agent's full transcript, so that the main agent can verify every step it took.
  \choice The sub-agent's tool schemas, so that the main agent could repeat the work itself.
  \choice Nothing --- the sub-agent writes its result to long-term memory, which the main agent reads later.
\end{choices}
\dontknow
\itemmeta{T11 --- Sub-agents and A2A}{understand}{medium}{A17}
\rationale{The point of delegation is context isolation; (b) would defeat it entirely.}

% ===== Observability and quality assurance ===========================

% --- Item 18: T12 / understand / easy ---------------------------------
\qitem Which set of measurements is most relevant for monitoring an agentic system in operation?
\begin{choices}
  \CorrectChoice Run throughput and latency, success and failure rates, model and token usage, tool usage, and permission decisions.
  \choice CPU temperature, disk fragmentation, and the number of open file descriptors.
  \choice The number of lines of code in the harness and its test coverage.
  \choice The model's parameter count, context window size, and training cut-off date.
\end{choices}
\dontknow
\itemmeta{T12 --- Observability and telemetry}{understand}{easy}{A18}
\rationale{(c) measures the codebase and (d) the model, neither of which changes while the system runs.}

% --- Item 19: T12 / apply-analyse / medium ----------------------------
\qitem Your dashboard shows that the tool failure rate rose from 2\% to 20\% this morning. What do you
need \emph{in addition} to the aggregate metrics in order to find the cause?
\begin{choices}
  \choice A longer retention period for the same aggregate counters.
  \choice A second dashboard showing the same metrics with a finer colour scale.
  \CorrectChoice Structured logs or traces that let you inspect what happened inside individual runs.
  \choice More frequent scraping, so that the same counters are sampled at a finer time resolution.
\end{choices}
\dontknow
\itemmeta{T12 --- Observability and telemetry}{apply-analyse}{medium}{A19}
\rationale{Aggregates detect; traces diagnose. (d) improves resolution of the same aggregate signal and
still cannot explain a single run.}

% --- Item 20: T13 / understand / medium -------------------------------
\qitem An end-to-end test instructs the agent to create \texttt{output/answer.md} containing the text
\texttt{forty-two}. Which assertion is robust against the agent's non-determinism?
\begin{choices}
  \choice The agent called \texttt{create\_folder} before \texttt{create\_file}.
  \CorrectChoice The file exists at that path and its contents are exactly the expected text.
  \choice The agent needed exactly two tool calls to complete the task.
  \choice The agent's final message contains the sentence ``I have created the file.''
\end{choices}
\dontknow
\itemmeta{T13 --- Quality assurance}{understand}{medium}{A20}
\rationale{Assert on resulting state, not on the path taken or the wording chosen --- both vary between
runs even when the outcome is correct.}

\end{questions}
